\documentclass[11pt,a4paper]{report}

\usepackage[margin=1in]{geometry}
\usepackage{listings}
\usepackage{xcolor}
\usepackage{booktabs}
\usepackage{longtable}
\usepackage{hyperref}
\usepackage{caption}
\usepackage{amsmath}
\usepackage{enumitem}
\usepackage{fancyhdr}

\hypersetup{colorlinks=true, linkcolor=blue, urlcolor=blue, citecolor=blue}

\definecolor{codebg}{rgb}{0.96,0.96,0.96}
\definecolor{codekw}{rgb}{0.0,0.0,0.6}
\definecolor{codecomment}{rgb}{0.3,0.5,0.3}
\definecolor{codestring}{rgb}{0.6,0.0,0.0}
\definecolor{okgreen}{rgb}{0.0,0.5,0.0}
\definecolor{badred}{rgb}{0.7,0.0,0.0}

\lstdefinestyle{cppstyle}{
  language=C++,
  backgroundcolor=\color{codebg},
  basicstyle=\ttfamily\small,
  commentstyle=\color{codecomment}\itshape,
  keywordstyle=\color{codekw}\bfseries,
  stringstyle=\color{codestring},
  numbers=left,
  numberstyle=\tiny\color{gray},
  numbersep=8pt,
  frame=single,
  rulecolor=\color{gray!40},
  breaklines=true,
  showstringspaces=false,
  tabsize=4,
  columns=fullflexible,
  keepspaces=true
}
\lstset{style=cppstyle}

\pagestyle{fancy}
\fancyhf{}
\fancyhead[L]{APARA C Compiler \& Engine --- Handoff Meeting Report}
\fancyhead[R]{\thepage}
\renewcommand{\headrulewidth}{0.4pt}

\newcommand{\filepath}[1]{\noindent\textbf{Location:}\ \texttt{#1}\par}

\title{\textbf{APARA C Compiler \& Engine Simulator \\ Handoff Meeting Report} \\
\large Complete compiler status, every engine source change with cause and
effect, exact file/line locations, and remaining known gaps}
\author{APARA C Compiler Project}
\date{\today}

\begin{document}

\maketitle
\tableofcontents

% ============================================================================
\chapter{Executive Summary}
% ============================================================================

\section{What is being delivered}

\begin{enumerate}[itemsep=2pt]
  \item \textbf{A feature-complete C compiler for the APARA VLIW ISA}
        (repo: \texttt{github.com/mohithkota/APARA\_COMPILER}), usable like
        gcc: clone, \texttt{export APARA\_TOOLS=...},
        \texttt{./apara-cc prog.c --run}.
  \item \textbf{A verification result}: 1{,}300+ C programs (hand-written
        golden suites + two randomized fuzz campaigns) compiled and executed
        on the simulator, with \emph{every} \texttt{results[]} slot matching
        native \texttt{gcc} \emph{bit-for-bit}; a coverage auditor proves all
        102 (instruction $\times$ sub-op $\times$ type) combinations the
        compiler emits are exercised. \textbf{Zero known compiler defects.}
  \item \textbf{Six simulator source fixes} (this report, Chapter~3) ---
        applied \emph{locally only}, shipped as a reviewable installer
        (\texttt{engine\_patches/apply\_engine\_fixes.py}); \textbf{nothing
        has been pushed to the upstream engine repository} --- upstream
        absorption is this meeting's decision.
  \item \textbf{A map of remaining engine gaps} (Chapter~4) --- reported,
        deliberately \emph{not} fixed; none affect compiler users.
\end{enumerate}

\section{Headline numbers}

\begin{center}
\begin{tabular}{@{}lr@{}}
\toprule
Golden regression gate (7 suites) & 71/71 PASS \\
Floating point fp01--fp09 (bit-exact vs gcc) & 50/50 checks \\
Fuzz campaign v1 (whole-program shapes) & 275/275 PASS \\
fuzz1000 campaign (1{,}000 seeds + 13 directed) & 613 PASS, 0 FAIL, 0 HANG \\
ISA coverage audit (instr $\times$ sub-op $\times$ type) & 102/102 \\
Real bugs found \& fixed during verification & 10 compiler-side + 6 simulator \\
\bottomrule
\end{tabular}
\end{center}

% ============================================================================
\chapter{Compiler Status}
% ============================================================================

\section{Feature completeness}

All features are verified against \texttt{gcc} on the simulator.

\begin{longtable}{@{}p{8.2cm}p{5.6cm}@{}}
\toprule
Feature & Evidence \\
\midrule
Full integer C: all widths signed/unsigned, all operators, all control flow,
recursion & \texttt{feature\_sweep} 16/16, \texttt{universal} 21/21 \\
Structs (nested), unions, bit-fields, enums, typedefs, statics, goto,
designated initializers & \texttt{feature\_sweep}, \texttt{fuzz1000/d10} \\
Pointers: arithmetic, deref, \texttt{\&x}, \texttt{p[i]}, comparison,
\texttt{q-p}, struct pointers & \texttt{pointer\_bugs} 15/15 \\
Arrays 1D/2D/3D, global+local, initialized & \texttt{fuzz1000/d10} \\
Floating point f32+f64: arithmetic, comparisons, casts, params, arrays,
\texttt{sqrt} & fp01--09 = 50/50, \emph{bit-exact} \\
Function pointers (callbacks, dispatch tables, \texttt{\&f}, \texttt{==},
returns) --- via a compile-time linker pass; no assembler change needed &
\texttt{fnptr} 16/16 \\
Variadic functions; $>$4 named arguments (unified stack-passing convention)
& \texttt{vararg} 12/12, \texttt{many\_args} 12/12 \\
Vector intrinsics: \texttt{\$v} +/-/* at all 6 widths, \texttt{\$dot}
($\pm$accumulate), \texttt{\$vreduce} +/max, \texttt{\$cmov} (6 conds),
\texttt{\$slice}, \texttt{\$pack}, \texttt{\$fsqrt} & \texttt{fuzz1000}
d03--d08 \\
Optimizer: 28-register dynamic allocation + spilling, CSE, LICM,
loop-carried register promotion, VLIW list scheduling \& bundling (with the
hardware's 4-ld/st and 1-div/sqrt lane limits enforced) & up to $-74\%$
static bundles \\
\bottomrule
\end{longtable}

Out of scope by design: real string handling (not needed for the
accelerator); NaN semantics in float comparisons; unsigned 64-bit ordered
comparison of values $\ge 2^{63}$ (documented).

\section{How correctness is established}

Every test follows the \texttt{results[]} convention. The compiler compiles
the \emph{same source} natively with gcc and records gcc's answers as the
expected DMEM post-conditions; the simulator then checks every slot. The
verdict is therefore ``\emph{produced the same values as gcc}'', never
``looks right''. The fuzz generators are UB-free by construction (bounded
loops/recursion, masked shifts, guarded divisors, forced 64-bit arithmetic
width, no NaN), so any mismatch is a real toolchain defect --- this is
exactly how the 16 bugs below were found.

% ============================================================================
\chapter{Engine Source Changes: Cause, Effect, Patch, Result}
% ============================================================================

Base directory for all paths in this chapter:

\begin{center}
\texttt{engine\_new/AjitHpcAccelRepo/AjitHpcAccel/engine\_isp/assembler/src/}
\end{center}

Line numbers are for the \emph{patched} tree on the development machine (the
same code the installer produces); in the pristine upstream tree, locate by
the function name --- each section quotes the exact original text to search
for. Changes marked \textbf{RE-APPLIED} were first made during the July
17--18 floating-point bring-up, deployed as binaries only, and later lost to
a rebuild (Section~\ref{sec:incident}); they are now permanent in source.

% ----------------------------------------------------------------------------
\section{Fix 1: \texttt{\$fsqrt} was an unimplemented stub (NEW)}

\filepath{MachineRun.cpp, line 599 --- McodeFsqrtInstruction::Execute}

\textbf{Cause.} The \texttt{Execute} method was a placeholder that printed
``\texttt{McodeFsqrtInstruction yet to be implemented}'' and wrote 0.

\textbf{Effect.} Every \texttt{\$fsqrt} returned 0:
\texttt{(long long)sqrt(4.0)} $\to$ 0 instead of 2. Anyone validating
hardware square-root against this simulator would have compared against
zeros. (Found by directed test \texttt{fuzz1000/directed/d08}.)

\textbf{Original code (search for this in the pristine tree):}
\begin{lstlisting}
void McodeFsqrtInstruction::Execute     (McodeMachine* mc)
{
    McodeRoot::Error("McodeFsqrtInstruction yet to be implemented", this);
    this->McodeInstruction::Set_Out_Arg (0, 0);
    return;
}
\end{lstlisting}

\textbf{Patch applied} (routes through the existing, already-correct ALU
dispatch --- \texttt{\_\_alu\_operation\_\_} already implements
\texttt{\_\_FSQRT} via \texttt{fp\_sqrt}; no new arithmetic was written):
\begin{lstlisting}
void McodeFsqrtInstruction::Execute     (McodeMachine* mc)
{
    uint64_t ovalue;
    ___execute_alu_operation___ (this->McodeInstruction::Get_Opcode(),
            this->Get_Dest_Type(),
            this->Get_Src_Type(),
            this->McodeInstruction::Get_In_Arg(0),
            0, 0, 0,
            ovalue);
    this->McodeInstruction::Set_Out_Arg (0, ovalue);
    if(__global_verbose_flag)
        McodeRoot::Info(Int64ToString (this->McodeInstruction::Get_Address()) + ": Fsqrt!");
    return;
}
\end{lstlisting}

\textbf{After the patch.} Simulator run-log shows
\texttt{0x4010000000000000 $\to$ 0x4000000000000000} ($\sqrt{4.0}=2.0$,
IEEE-exact); d08 = 8/8 vs gcc for \texttt{sqrt}/\texttt{sqrtf} over
f64 and f32.

% ----------------------------------------------------------------------------
\section{Fix 2: float casts hit an error branch, output left uninitialized (RE-APPLIED)}

\filepath{McodeExecute.cpp, lines 27--33 --- \_\_\_execute\_cast\_operation\_\_\_}

\textbf{Cause.} The complete cast implementation was guarded by
\texttt{if(!dest\_type.Get\_Float\_Flag() \&\& !src\_type.Get\_Float\_Flag())};
any float-typed cast fell into an \texttt{else} that printed
``\texttt{Float conversions in cast not supported as yet.}'' and returned
\emph{without writing the output array}.

\textbf{Effect.} The destination register received whatever happened to be
in host memory --- observed values like \texttt{0x7ffd83e474a0} (a host
stack address). Nondeterministic garbage in every float$\leftrightarrow$int
conversion.

\textbf{Original code (search):}
\begin{lstlisting}
    if(!dest_type.Get_Float_Flag()
            && !src_type.Get_Float_Flag())
    { /* full implementation */ }
    else
    {
        McodeRoot::Error ("Float conversions in cast not supported as yet.", NULL);
    }
\end{lstlisting}

\textbf{Patch applied.} Guard replaced with \texttt{if(1)} and the error
branch removed --- the body it guards forwards to
\texttt{\_\_\_cast\_operation\_\_\_}, which dispatches
int$\leftrightarrow$float itself (Fix~3), so every cast can take the one
real path.

\textbf{After the patch.} fp01--fp09 50/50 bit-exact (together with Fixes
3, 4, 6); float casts appear in 600+ passing fuzz programs.

% ----------------------------------------------------------------------------
\section{Fix 3: int$\leftrightarrow$float dispatch called the two helpers swapped (RE-APPLIED)}

\filepath{McodeOperations.cpp, lines 330--343 --- \_\_\_cast\_operation\_\_\_}

\textbf{Cause.} A transposition: the integer-source case called
\texttt{cast\_float\_to\_int} and the float-source/int-dest case called
\texttt{cast\_int\_to\_float}.

\textbf{Effect.} Every conversion read its operand as the wrong
representation (integer bits interpreted as IEEE and vice versa) ---
garbage in both directions.

\textbf{Original code (search):}
\begin{lstlisting}
            if(!src_type.Get_Float_Flag())
                w = cast_float_to_int (src_type, dest_type, v);
            else if(!dest_type.Get_Float_Flag())
                w = cast_int_to_float (src_type, dest_type, v);
\end{lstlisting}

\textbf{Patch applied} (the two call sites exchanged):
\begin{lstlisting}
            if(!src_type.Get_Float_Flag())
                w = cast_int_to_float (src_type, dest_type, v);
            else if(!dest_type.Get_Float_Flag())
                w = cast_float_to_int (src_type, dest_type, v);
\end{lstlisting}

\textbf{After the patch.} See Fix 2 (all casts flow through here).

% ----------------------------------------------------------------------------
\section{Fix 4: scalar 32-bit casts misclassified as vector casts (NEW)}

\filepath{McodeOperations.cpp, lines 314--321 --- \_\_\_cast\_operation\_\_\_}

\textbf{Cause.} Vector-ness was inferred from the element count,
\texttt{is\_vector\_cast = (in\_vals.size() > 1)}. \texttt{Break\_Vector}
splits even a \emph{scalar} 32-bit source register into two entries, so
every scalar cast with a 32-bit side was treated as a vector cast --- and
vector casts mask each output to the destination lane width.

\textbf{Effect.} \texttt{(int)(-6.0f)} stored
\texttt{0x00000000FFFFFFFA} instead of \texttt{0xFFFFFFFFFFFFFFFA}: the
conversion helper's correct sign-extended result was truncated back to 32
bits. Every negative float$\to$int result stored into a 64-bit word was
wrong. (Caught by regression test fp03.)

\textbf{Original code (search):}
\begin{lstlisting}
    int is_vector_cast = (in_vals.size() > 1);
\end{lstlisting}

\textbf{Patch applied} (trust the instruction's declared types):
\begin{lstlisting}
    int is_vector_cast = (src_type.Get_Vector_Flag() && dest_type.Get_Vector_Flag());
\end{lstlisting}

\textbf{After the patch.} fp03 4/4 restored; gate 61/61; true vector casts
behave exactly as before (both flags set).

% ----------------------------------------------------------------------------
\section{Fix 5: unsigned \texttt{\$vreduce} sign-extended its lanes (NEW)}

\filepath{McodeOperations.cpp, lines 150--156 --- \_\_vreduce\_operation\_\_}

\textbf{Cause.} A shared lane value \texttt{r} was computed with
\texttt{Sign\_Extend\_64} \emph{before} the signed/unsigned branch. The
signed branch shadows it with its own copy, but the \emph{unsigned} branch
consumed the outer, sign-extended value.

\textbf{Effect.} \texttt{\_\_vreduce\_vu8(0x01f2030405060708)} returned
\texttt{0x14} (20) instead of \texttt{0x114} (276) --- the \texttt{0xf2}
lane entered the sum as $-14$ instead of $242$;
\texttt{\_\_vreduce\_max\_vu*} picked wrong maxima. In effect, the entire
unsigned half of \texttt{\$vreduce} behaved as signed. (Found by directed
test d06.)

\textbf{Original code (search):}
\begin{lstlisting}
            // sign extend to 64-bits.
            int64_t r = (int64_t) Sign_Extend_64(src_type.Get_Nbits()-1, ele);
            if(signed_flag)
            {
\end{lstlisting}

\textbf{Patch applied} (\texttt{ele} arrives already lane-masked from
\texttt{Break\_Vector}, i.e.\ zero-extended; signed branch untouched):
\begin{lstlisting}
            uint64_t r = ele;
            if(signed_flag)
            {
                int64_t r = (int64_t) Sign_Extend_64(...);  // as before
\end{lstlisting}

\textbf{After the patch.} d06 = 8/8: \texttt{+} and \texttt{\$max}
reductions at \texttt{vi8/vu8/vi16/vu16/vi32/vu32} all match gcc.

% ----------------------------------------------------------------------------
\section{Fix 6: \texttt{cast\_int\_to\_float} ternary promotion of negatives (RE-APPLIED)}

\filepath{McodeFpuUtils.cpp, lines 437--450 --- cast\_int\_to\_float}

\textbf{Cause.} C's conditional operator promotes both arms to a common
type. In
\texttt{cond ? (u \& 0xffffffff) : (int) u}
the unsigned arm's type is \texttt{uint64\_t}, so the signed arm's negative
\texttt{int} was converted back to a huge unsigned value before becoming a
\texttt{double}.

\textbf{Effect.} \texttt{(float)(-3)} became $2^{64}-3 \approx
1.8\times10^{19}$; every \emph{negative} int$\to$float cast was wrong.
(All positive casts were fine --- which is how it escaped initial testing.)

\textbf{Original code (search):}
\begin{lstlisting}
    double x = (src_type.Get_Unsigned_Flag() ?  (u & 0xffffffff) : (int) u);
\end{lstlisting}

\textbf{Patch applied} (explicit branch; sign-extend from the \emph{actual}
source width --- this also fixes the old hard-coded 32-bit mask on the
unsigned arm):
\begin{lstlisting}
    double x;
    if(src_type.Get_Unsigned_Flag())
    {
        x = (double) (u & __mmask__(snbits));
    }
    else
    {
        int64_t sv = ((int64_t) (u << (64 - snbits))) >> (64 - snbits);
        x = (double) sv;
    }
\end{lstlisting}

\textbf{After the patch.} fp03 (positive \emph{and} negative int$\to$float,
$\times$, $\div$, $+$) = 4/4 vs gcc.

% ----------------------------------------------------------------------------
\section{The binary-only-fix incident, and the safeguards}
\label{sec:incident}

Fixes 2, 3 and 6 were first made during the July 17--18 floating-point
bring-up, but deployed as \emph{binaries only}; the source files stayed
pristine. On July 19 a rebuild (for Fix 1) regenerated all binaries from
the unfixed source and silently reverted them. The regression gate caught
it immediately (fp03 failing, 60/61), and everything was re-applied
\emph{in source}. Safeguards now in place:

\begin{itemize}[itemsep=2pt]
  \item \texttt{engine\_patches/apply\_engine\_fixes.py --check <src>}
        verifies all six fixes are present (exact-text match); \texttt{--apply}
        installs them idempotently on a pristine tree and reports a CONFLICT
        if upstream has diverged.
  \item The full gate is re-run after \emph{every} toolchain rebuild.
\end{itemize}

% ============================================================================
\chapter{Remaining Gaps in the Engine (Reported, NOT Fixed)}
% ============================================================================

These were found during the audit and deliberately left untouched ---
they are outside what the compiler emits, so compiler users are unaffected
(the 102/102 coverage audit is the proof: everything the compiler emits is
exercised and correct).

\begin{longtable}{@{}p{5.2cm}p{4.2cm}p{4.6cm}@{}}
\toprule
Gap \& exact location & Failure mode & Impact \\
\midrule
\textbf{\texttt{\$vreduce *} always returns 0} ---
\texttt{McodeOperations.cpp:108} (author's own \texttt{// TODO: add
mreduce}) and \texttt{:161}: \texttt{s\_result = s\_result * r} folds from
an accumulator initialized to \textbf{0} (multiplicative identity should
be 1). Validity check \emph{allows} it for 32-bit lanes. &
Silent wrong value (always 0) for hand-written \texttt{\$vreduce *} &
None for compiler users (compiler emits only \texttt{+}/\texttt{\$max});
\textbf{newly found in this audit} \\
\addlinespace
f4/f8/f16 float widths: \texttt{fp\_div} / \texttt{fp\_sqrt} and cast
helpers end in \texttt{default: assert(0)} ---
\texttt{McodeFpuUtils.cpp:407, 426, 460} & Loud crash & None ---
compiler emits f32/f64 only \\
\addlinespace
\texttt{\$vreduce} sub-op restrictions: ADD (sub-64 lanes), MUL (32-bit
only), MAX/MIN; AND/OR/XOR/XNOR rejected ---
\texttt{McodeInstructions.cpp:620--655} & Loud reject & None \\
\addlinespace
Native \texttt{\$abs}/\texttt{\$max}/\texttt{\$min} instructions
grammar-legal but broken in the assembler/simulator (June finding) &
Wrong values & None --- compiler lowers \texttt{\_\_abs/\_\_max/\_\_min}
to \texttt{\$cmov} \\
\addlinespace
\texttt{\$pack} validity (\texttt{packed\_nbits \% word\_nbits == 0},
\texttt{McodeClasses.hpp:412}) is checked \emph{silently}; the assembler
reports only ``there were invalid instructions'' & Hard-to-diagnose
assembly failure & Minor developer-experience issue \\
\addlinespace
\texttt{mcode\_align} prints lane-placement errors ($>$4 ld/st or $>$1
div/sqrt per bundle) to stderr but \emph{still emits} the bundle with the
CTI not in lane 0 (\texttt{McodeProgram.cpp},
\texttt{alignFullBundleToLanes} / \texttt{Print}) --- downstream this sent
a \texttt{\$return} into the middle of a bundle & Undefined control flow
from a ``successful'' build & Compiler's bundler now enforces the limits
itself; recommend align should hard-fail \\
\addlinespace
IMEM is $\sim$0x800 words ($\approx$ 240 bundles), half the documented
size (June finding) & Large programs don't fit & The one real capacity
limit; compiler reports it \\
\bottomrule
\end{longtable}

% ============================================================================
\chapter{Verification Evidence \& How to Reproduce}
% ============================================================================

\section{What was run}

\begin{itemize}[itemsep=2pt]
  \item \textbf{Golden gate} (71 hand-written tests, 7 suites):
        \texttt{bash testing/run\_gate.sh} $\to$ 71/71.
  \item \textbf{Fuzz v1} (\texttt{testing/fuzz/}): 400 seeded random
        whole-programs $\to$ 275 executed, 275 PASS.
  \item \textbf{fuzz1000} (\texttt{testing/fuzz1000/}): 13 directed + 1000
        seeded programs $\to$ 613 executed, 613 PASS, 0 FAIL, 0 HANG;
        coverage audit 102/102 (\texttt{cov\_scan.py --report}).
  \item Every executed program's \texttt{results[]} compared slot-by-slot,
        bit-for-bit, against native gcc on the same source.
\end{itemize}

\section{Reproduce on any machine}

\begin{lstlisting}[language=bash]
git clone git@github.com:mohithkota/APARA_COMPILER.git && cd APARA_COMPILER
pip install pycparser
export APARA_TOOLS=/path/to/engine_isp/assembler/bin

# verify (or, after approval, install) the six simulator fixes:
python3 engine_patches/apply_engine_fixes.py --check /path/to/engine_isp/assembler/src

./apara-cc mytest.c --run          # gcc-style compile + execute + verdict
bash testing/run_gate.sh           # full 71-test gate
bash testing/fuzz1000/run_all.sh 1 1000   # the whole campaign, if desired
\end{lstlisting}

\section{Bugs found in the compiler itself during verification}

For completeness (all fixed, all documented in
\texttt{compiler/STATUS.md} with dates): LICM cross-function alias-map
collision (infinite loop); struct-pointer subscript reading DMEM address 0;
bundler exceeding the hardware's 4-ld/st / 1-div-sqrt lanes; literal-suffix
(\texttt{LL}/\texttt{u}) global initializers zeroed; struct-global float
field initializers zeroed; 3D-array init/access stride mismatch; missing
float$\to$int cast on \texttt{sqrt} results; unsigned 64-bit
\texttt{/ \% >>} on the signed path; \texttt{\_\_nop()} silently deleted
by the bundler; fp signature bug in the fuzz generator. Ten real defects,
every one found by differential testing against gcc.

% ============================================================================
\chapter{Decisions Requested in This Meeting}
% ============================================================================

\begin{enumerate}[itemsep=4pt]
  \item \textbf{Approve the six simulator fixes} (Chapter~3) for the
        upstream tree. Each is minimal and local; the installer
        (\texttt{apply\_engine\_fixes.py}) applies them exactly, or they can
        be re-typed by hand from this report. \emph{Until approved, upstream
        is untouched.}
  \item \textbf{Acknowledge the remaining gaps} (Chapter~4) --- especially
        the newly-found \texttt{\$vreduce *} silent-zero and the
        recommendation that \texttt{mcode\_align} hard-fail on lane errors.
  \item \textbf{Hardware validation ownership}: the compiler author is
        off-campus; the repo is self-contained (clone $\to$ gate $\to$
        hardware runs), so any lab member can execute the hardware
        confirmation of the post-June features (FP, function pointers,
        variadics, $>$4 args).
  \item \textbf{IMEM size}: confirm whether the half-size IMEM is
        intentional; it is the only capacity limit users will hit.
\end{enumerate}

\end{document}
